
% ---
% Dedicatória (opcional)
% ---
\begin{dedicatoria}
   \vspace*{\fill}
   \vspace*{\fill}
   \vspace*{\fill}
   \vspace*{\fill}
   \flushright
   \noindent
   \textit{À alguém...}
   \vspace*{48pt}
\end{dedicatoria}
% ---

% ---
% Agradecimentos (opcional)
% ---
\begin{agradecimentos}
Liste a quem deseja agradecer (opcional) em ordem decrescente de importância, evitando um número muito extenso de pessoas.
Os agradecimentos parágrafo parágrafo parágrafo parágrafo parágrafo parágrafo parágrafo parágrafo parágrafo parágrafo parágrafo parágrafo parágrafo parágrafo parágrafo parágrafo parágrafo parágrafo parágrafo parágrafo parágrafo parágrafo parágrafo parágrafo parágrafo parágrafo parágrafo parágrafo parágrafo parágrafo parágrafo parágrafo parágrafo parágrafo parágrafo parágrafo parágrafo parágrafo parágrafo parágrafo parágrafo 

parágrafo parágrafo parágrafo parágrafo parágrafo parágrafo parágrafo parágrafo parágrafo parágrafo parágrafo parágrafo parágrafo parágrafo parágrafo parágrafo parágrafo parágrafo parágrafo parágrafo parágrafo parágrafo parágrafo parágrafo parágrafo parágrafo parágrafo parágrafo parágrafo parágrafo parágrafo parágrafo parágrafo parágrafo parágrafo parágrafo parágrafo parágrafo parágrafo parágrafo parágrafo parágrafo parágrafo parágrafo parágrafo parágrafo parágrafo parágrafo parágrafo parágrafo parágrafo parágrafo parágrafo parágrafo parágrafo 

\end{agradecimentos}
% ---

% ---
% Epígrafe (opcional)
% ---
\begin{epigrafe}
    \vspace*{\fill}
	\begin{flushright}
		''Elemento opcional, que traz a citação de um pensamento seguido de autoria e que represente a gênese da obra (seu trabalho). Não há necessidade de escrever o título “Epígrafe”. Usar letras minúsculas alinhadas à direita da parte inferior da folha'' (Alguém, Livro, ano)
        \vspace*{48pt}
	\end{flushright}
\end{epigrafe}